% =========================================================
% Proof: Duality Principle
% =========================================================

\newpage
\phantomsection
\label{prf:duality-principle}

\begin{remark*}[Return]
\hyperref[prop:duality-principle]{Return to Proposition~\ref*{prop:duality-principle}.}
\end{remark*}

\begin{theorem*}[Duality Principle]
Let $\Phi$ be a statement about partially ordered sets expressed in terms
of $\leq$ that is true in every partially ordered set.  Then the dual
statement $\Phi^{\mathrm{op}}$, obtained by reversing the order, is also
true in every partially ordered set.
\end{theorem*}

\begin{proof}[Proof]
Let $P$ be any partially ordered set.  The dual $P^{\mathrm{op}}$ is
also a partially ordered set
(\hyperref[def:dual-ordered-set]{Definition~\ref*{def:dual-ordered-set}}).
Since $\Phi$ holds in every partially ordered set, $\Phi$ holds in
$P^{\mathrm{op}}$.  The statement $\Phi$ applied to $P^{\mathrm{op}}$ is,
by definition of $P^{\mathrm{op}}$, precisely the statement $\Phi^{\mathrm{op}}$
applied to $P$.  Therefore $\Phi^{\mathrm{op}}$ holds in $P$.  Since $P$
was arbitrary, $\Phi^{\mathrm{op}}$ holds in every partially ordered set.
\end{proof}

\begin{remark*}[Proof structure]
A one-step argument: the dual $P^{\mathrm{op}}$ is a valid poset, so
any universal statement about posets applies to it.  Reading that
statement back in the original order of $P$ yields the dual statement.
\end{remark*}

\begin{remark*}[Dependencies]
\hyperref[def:dual-ordered-set]{Dual ordered set},
\hyperref[def:lattice-order-partial-order]{partial order}.
\end{remark*}
